\documentclass[11pt,letterpaper]{article}
\usepackage{cornell-notes}
\usepackage{needspace}

\title{Chapter 17: Web Applications}
\author{}
\date{}

\setDocTitle{Chapter 17: Web Applications}
\setDocSubtitle{Application Security Review}
\setCornellCollection{Software Security Assessment}
\setCornellUnitType{Chapter}
\setCornellUnitNumber{17}
\setCornellUnitTitle{Web Applications}
\setCornellCompanionLabel{Study and audit reference}
\setCornellSourceRange{pp. 1007--1081}
\begin{document}
\maketitle

\begin{center}
\small\color{Meta}\textbf{Coverage range:} pp. 1007--1081 \quad\textbullet\quad \textbf{Format:} Cornell cue-and-notes
\end{center}
\vspace{0.25em}
\begin{CornellOverviewBox}{Review Orientation}
\textbf{Central problem.} Audit web applications across HTTP, client-controlled state, sessions, layered architecture, access control, encryption, navigation, and common injection paths.
\textbf{Why it matters.} Security reviewers use this material to connect attacker-controlled conditions to violated invariants, consequential operations, and defensible remediation.
\textbf{Prerequisites.} HTTP, cookies, server-side execution, databases, browser behavior, authentication, authorization, and input validation.
\textbf{Assessment relationship.} It combines earlier design, parsing, state, and operational techniques in the multi-tier web environment.
\end{CornellOverviewBox}
\section{Learning Objectives}
\begin{itemize}
\item Explain the security significance of static and dynamic mechanisms.
\item Identify attacker influence and failed assumptions associated with http versions and headers.
\item Trace security-relevant data or state through methods.
\item Evaluate whether controls addressing parameters and forms preserve the intended invariant.
\item Audit implementation and error paths related to client addresses and referrers.
\item Distinguish safe and unsafe uses of embedded state.
\item Diagnose a failure involving http authentication from concrete evidence.
\item Verify remediation and test coverage for cookies and sessions.
\end{itemize}
\section{Cornell Cue-and-Notes}
\Needspace{8\baselineskip}
\subsection{Web Technology and HTTP}
\begin{longtable}{>{\raggedright\arraybackslash\columncolor{CornellCueBackground}}p{0.285\textwidth}|>{\raggedright\arraybackslash}p{0.625\textwidth}}
\rowcolor{Navy}\color{white}\textbf{Cue / Recall Prompt} & \color{white}\textbf{Technical Notes and Audit Reasoning} \\
\endfirsthead
\rowcolor{Navy}\color{white}\textbf{Cue / Recall Prompt} & \color{white}\textbf{Technical Notes and Audit Reasoning} \\
\endhead
\term{Static and dynamic mechanisms}\\[-0.1em]{\small Which component executes a request?} & Static files, CGI programs, server APIs, includes, transformations, and scripts follow different mappings and privilege models. Extension and path handling can expose source or invoke unintended handlers.\par\smallskip\textbf{Audit move:} Map URL patterns and content types to exact handlers and execution identities. \\\hline
\term{HTTP versions and headers}\\[-0.1em]{\small Which request metadata influences security decisions?} & Hosts, content framing, forwarding information, content types, and caching directives can change routing and interpretation. Headers are attacker-controlled unless set and authenticated by a trusted intermediary.\par\smallskip\textbf{Audit move:} Inventory security-sensitive headers and establish their trusted source and normalization. \\\hline
\term{Methods}\\[-0.1em]{\small Does every method enforce the same policy?} & Read, write, delete, diagnostic, and extension methods may reach different handlers. Restricting only ordinary retrieval leaves alternate method paths exposed.\par\smallskip\textbf{Audit move:} Enumerate accepted methods per route and test authorization and body parsing for each. \\\hline
\term{Parameters and forms}\\[-0.1em]{\small Where can duplicate or missing parameters alter logic?} & Query, path, form, multipart, and JSON-like values can repeat, be empty, or use alternate encodings. Framework binding rules determine which value reaches business logic.\par\smallskip\textbf{Audit move:} Test duplicates, arrays, conflicting sources, encodings, and boundary lengths. \\\hline
\end{longtable}
\begin{CornellSummaryBox}{Section Summary}
Web applications expose server behavior through protocol fields, server mappings, dynamic execution, and client-controlled navigation.
\end{CornellSummaryBox}
\Needspace{8\baselineskip}
\subsection{State and Authentication}
\begin{longtable}{>{\raggedright\arraybackslash\columncolor{CornellCueBackground}}p{0.285\textwidth}|>{\raggedright\arraybackslash}p{0.625\textwidth}}
\rowcolor{Navy}\color{white}\textbf{Cue / Recall Prompt} & \color{white}\textbf{Technical Notes and Audit Reasoning} \\
\endfirsthead
\rowcolor{Navy}\color{white}\textbf{Cue / Recall Prompt} & \color{white}\textbf{Technical Notes and Audit Reasoning} \\
\endhead
\term{Client addresses and referrers}\\[-0.1em]{\small Can transport metadata authenticate a user?} & Addresses may be shared, proxied, spoofed in headers, or change over time, while referrer-like headers are optional and client-controlled. They are weak context signals, not primary identity proof.\par\smallskip\textbf{Audit move:} Do not authorize sensitive actions solely from address or referrer metadata. \\\hline
\term{Embedded state}\\[-0.1em]{\small Can hidden fields or URLs be trusted?} & State placed in forms or URLs returns under client control. It must be integrity-protected or rederived server-side, and secrets should not appear in leak-prone locations.\par\smallskip\textbf{Audit move:} Classify every client-returned state value and verify server-side validation or protection. \\\hline
\term{HTTP authentication}\\[-0.1em]{\small Where are credentials protected and bound?} & Authentication mechanisms differ in replay resistance, channel dependence, credential exposure, and proxy behavior. The application must bind the established identity to authorization and session state.\par\smallskip\textbf{Audit move:} Test downgrade, challenge handling, logout, credential transport, and identity propagation. \\\hline
\term{Cookies and sessions}\\[-0.1em]{\small What makes a session token secure?} & Tokens need sufficient unpredictability, protected transport, scoped storage, rotation, expiration, server-side invalidation, and resistance to fixation. Session data must not mix across users or concurrent transitions.\par\smallskip\textbf{Audit move:} Review generation, flags, issuance, rotation, lookup, logout, timeout, and revocation. \\\hline
\end{longtable}
\begin{CornellSummaryBox}{Section Summary}
Because HTTP is stateless, applications create state tokens whose integrity, confidentiality, lifetime, and binding become security-critical.
\end{CornellSummaryBox}
\Needspace{8\baselineskip}
\subsection{Architecture and Problem Areas}
\begin{longtable}{>{\raggedright\arraybackslash\columncolor{CornellCueBackground}}p{0.285\textwidth}|>{\raggedright\arraybackslash}p{0.625\textwidth}}
\rowcolor{Navy}\color{white}\textbf{Cue / Recall Prompt} & \color{white}\textbf{Technical Notes and Audit Reasoning} \\
\endfirsthead
\rowcolor{Navy}\color{white}\textbf{Cue / Recall Prompt} & \color{white}\textbf{Technical Notes and Audit Reasoning} \\
\endhead
\term{N-tier architecture}\\[-0.1em]{\small Where should validation and policy live?} & Presentation, business, and data tiers have different responsibilities. Security decisions must remain enforced at trusted service boundaries even when clients or presentation code appear to constrain input.\par\smallskip\textbf{Audit move:} Trace each sensitive action to the authoritative business-layer decision and data operation. \\\hline
\term{Client visibility and control}\\[-0.1em]{\small What can a browser user modify?} & Scripts, markup, hidden fields, requests, cookies, and navigation are visible and mutable. Client checks improve usability but cannot enforce server policy.\par\smallskip\textbf{Audit move:} Repeat every sensitive request without client code and with modified state. \\\hline
\term{Page flow}\\[-0.1em]{\small Can users invoke steps out of order?} & Links and wizard interfaces do not prevent direct requests, replay, or skipping. The server must validate prerequisites and state transitions for every action.\par\smallskip\textbf{Audit move:} Model the server-side state machine and test reordered, repeated, and omitted steps. \\\hline
\term{Authentication and authorization}\\[-0.1em]{\small Where do identity and object access diverge?} & A valid session proves a user identity, not permission for every object or function. Object identifiers, roles, tenants, and workflow state all require server-side checks.\par\smallskip\textbf{Audit move:} Build a role-object-action matrix and test horizontal and vertical access. \\\hline
\term{Encryption and impersonation}\\[-0.1em]{\small What does TLS protect, and what remains?} & Transport protection secures a connection when endpoints are validated, but does not repair application authorization, phishing, weak recovery, or compromised client state. Mixed or downgraded flows can expose tokens.\par\smallskip\textbf{Audit move:} Require protected transport for credentials and sessions and verify endpoint and redirect behavior. \\\hline
\end{longtable}
\begin{CornellSummaryBox}{Section Summary}
Layering clarifies responsibilities but creates cross-tier trust, serialization, and policy-consistency risks.
\end{CornellSummaryBox}
\Needspace{8\baselineskip}
\subsection{Common Vulnerabilities}
\begin{longtable}{>{\raggedright\arraybackslash\columncolor{CornellCueBackground}}p{0.285\textwidth}|>{\raggedright\arraybackslash}p{0.625\textwidth}}
\rowcolor{Navy}\color{white}\textbf{Cue / Recall Prompt} & \color{white}\textbf{Technical Notes and Audit Reasoning} \\
\endfirsthead
\rowcolor{Navy}\color{white}\textbf{Cue / Recall Prompt} & \color{white}\textbf{Technical Notes and Audit Reasoning} \\
\endhead
\term{SQL injection}\\[-0.1em]{\small How does input alter a database statement?} & Concatenating data into SQL allows quotes, operators, comments, or structure to change the intended query. Parameter binding separates values from grammar; authorization must still limit the resulting operation.\par\smallskip\textbf{Audit move:} Trace all database calls and require parameterized statements with least-privilege accounts. \\\hline
\term{OS and file interaction}\\[-0.1em]{\small Can web input select a command or filesystem object?} & Shell invocation, path traversal, temporary files, uploads, and unsafe permissions bridge remote input to operating-system authority. Web-server privilege amplifies impact.\par\smallskip\textbf{Audit move:} Prefer direct APIs, canonicalize paths, isolate uploads, and minimize process permissions. \\\hline
\term{XML and XPath injection}\\[-0.1em]{\small Where does structured input become query or document syntax?} & Concatenated values can alter XPath expressions, while XML parsers may process expensive, external, or ambiguous structures. Parser configuration and parameterization both matter.\par\smallskip\textbf{Audit move:} Use safe query construction and disable unnecessary external or expansive parser features. \\\hline
\term{Cross-site scripting}\\[-0.1em]{\small How does server data become browser code?} & Untrusted data inserted into HTML, attributes, URLs, styles, or script contexts can execute with the site's origin. Encoding must match the exact output context and occur at the final sink.\par\smallskip\textbf{Audit move:} Classify every output context and use context-specific encoding plus safe templating. \\\hline
\term{Threading and native-code issues}\\[-0.1em]{\small Can web concurrency or unsafe components violate isolation?} & Shared application state races across requests, and native extensions inherit memory and language hazards. High request concurrency makes rare timing flaws reachable.\par\smallskip\textbf{Audit move:} Audit shared mutable state and apply native-code review to modules handling web input. \\\hline
\end{longtable}
\begin{CornellSummaryBox}{Section Summary}
Web vulnerabilities occur where untrusted values cross into interpreters, files, shared state, or browser execution contexts.
\end{CornellSummaryBox}
\Needspace{8\baselineskip}
\subsection{Auditing Strategy}
\begin{longtable}{>{\raggedright\arraybackslash\columncolor{CornellCueBackground}}p{0.285\textwidth}|>{\raggedright\arraybackslash}p{0.625\textwidth}}
\rowcolor{Navy}\color{white}\textbf{Cue / Recall Prompt} & \color{white}\textbf{Technical Notes and Audit Reasoning} \\
\endfirsthead
\rowcolor{Navy}\color{white}\textbf{Cue / Recall Prompt} & \color{white}\textbf{Technical Notes and Audit Reasoning} \\
\endhead
\term{Route and entry-point inventory}\\[-0.1em]{\small What is the complete callable surface?} & Include documented and hidden routes, methods, APIs, files, handlers, administrative functions, callbacks, and legacy endpoints. Client navigation is not a complete inventory.\par\smallskip\textbf{Audit move:} Derive routes from code and configuration, then reconcile with crawling and traffic. \\\hline
\term{End-to-end data flow}\\[-0.1em]{\small How is one request transformed across tiers?} & Trace decoding, binding, validation, authorization, business logic, persistence, output encoding, and logging. A value can cross several grammars and identities in one request.\par\smallskip\textbf{Audit move:} Record each transformation and the guarantee required by the next consumer. \\\hline
\end{longtable}
\begin{CornellSummaryBox}{Section Summary}
A web review combines route inventory, trust-boundary analysis, session and authorization testing, sink tracing, and deployment inspection.
\end{CornellSummaryBox}
\begin{CornellLegacyBox}{Historical Context}
Specific server mechanisms and browser behavior evolve. Apply the enduring trust, state, output-context, and tier-boundary analysis using current framework and browser security guidance.
\end{CornellLegacyBox}
\section{Security-Review Checklist}
\begin{CornellChecklist}
\item \textbf{Scoping}
Map URL patterns and content types to exact handlers and execution identities.
\item \textbf{Attack-surface identification}
Inventory security-sensitive headers and establish their trusted source and normalization.
\item \textbf{Input and data-flow analysis}
Enumerate accepted methods per route and test authorization and body parsing for each.
\item \textbf{Trust-boundary review}
Test duplicates, arrays, conflicting sources, encodings, and boundary lengths.
\item \textbf{Candidate-point inspection}
Do not authorize sensitive actions solely from address or referrer metadata.
\item \textbf{Control-flow review}
Classify every client-returned state value and verify server-side validation or protection.
\item \textbf{Resource and lifetime review}
Test downgrade, challenge handling, logout, credential transport, and identity propagation.
\item \textbf{Error-path analysis}
Review generation, flags, issuance, rotation, lookup, logout, timeout, and revocation.
\item \textbf{Mitigation validation}
Trace each sensitive action to the authoritative business-layer decision and data operation.
\item \textbf{Evidence and reporting}
Repeat every sensitive request without client code and with modified state.
\item \textbf{Scoping}
Model the server-side state machine and test reordered, repeated, and omitted steps.
\item \textbf{Attack-surface identification}
Build a role-object-action matrix and test horizontal and vertical access.
\end{CornellChecklist}
\section{Chapter Synthesis}
Web application review treats the browser and all request data as untrusted, reconstructs server-side state and policy, and follows values across multiple tiers and grammar contexts. The highest-risk findings combine weak authorization or session binding with interpreter, filesystem, or operating-system access. Coverage must come from code and configuration as well as visible navigation.
\section{Self-Test Questions}
\begin{enumerate}
\item Which component executes a request?
\item Which request metadata influences security decisions?
\item Does every method enforce the same policy?
\item Where can duplicate or missing parameters alter logic?
\item Can transport metadata authenticate a user?
\item Can hidden fields or URLs be trusted?
\item \textbf{Scenario:} A multi-step transfer page hides the confirmation endpoint until earlier steps finish, but the endpoint accepts a direct request. What failed?
\item \textbf{Scenario:} A value is HTML-escaped and then inserted into a JavaScript string. Is the encoding necessarily correct?
\end{enumerate}
\clearpage
\subsection*{Answer Key}
\begin{enumerate}
\item Static files, CGI programs, server APIs, includes, transformations, and scripts follow different mappings and privilege models. Extension and path handling can expose source or invoke unintended handlers.
\item Hosts, content framing, forwarding information, content types, and caching directives can change routing and interpretation. Headers are attacker-controlled unless set and authenticated by a trusted intermediary.
\item Read, write, delete, diagnostic, and extension methods may reach different handlers. Restricting only ordinary retrieval leaves alternate method paths exposed.
\item Query, path, form, multipart, and JSON-like values can repeat, be empty, or use alternate encodings. Framework binding rules determine which value reaches business logic.
\item Addresses may be shared, proxied, spoofed in headers, or change over time, while referrer-like headers are optional and client-controlled. They are weak context signals, not primary identity proof.
\item State placed in forms or URLs returns under client control. It must be integrity-protected or rederived server-side, and secrets should not appear in leak-prone locations.
\item Client page flow was mistaken for server enforcement. The server must validate authenticated session, authorization, transaction state, integrity-protected parameters, and replay conditions on the confirmation request.
\item No. HTML and JavaScript string contexts have different grammars. Use context-specific encoding or safe serialization at the final sink and avoid constructing executable code from data.
\end{enumerate}
\section{Key-Term Recap}
\begin{longtable}{>{\raggedright\arraybackslash}p{0.27\textwidth}p{0.65\textwidth}}
\toprule
\textbf{Term} & \textbf{Working meaning for review} \\
\midrule
\endfirsthead
\toprule
\textbf{Term} & \textbf{Working meaning for review} \\
\midrule
\endhead
Static and dynamic mechanisms & Static files, CGI programs, server APIs, includes, transformations, and scripts follow different mappings and privilege models. \\
HTTP versions and headers & Hosts, content framing, forwarding information, content types, and caching directives can change routing and interpretation. \\
Methods & Read, write, delete, diagnostic, and extension methods may reach different handlers. \\
Parameters and forms & Query, path, form, multipart, and JSON-like values can repeat, be empty, or use alternate encodings. \\
Client addresses and referrers & Addresses may be shared, proxied, spoofed in headers, or change over time, while referrer-like headers are optional and client-controlled. \\
Embedded state & State placed in forms or URLs returns under client control. \\
HTTP authentication & Authentication mechanisms differ in replay resistance, channel dependence, credential exposure, and proxy behavior. \\
Cookies and sessions & Tokens need sufficient unpredictability, protected transport, scoped storage, rotation, expiration, server-side invalidation, and resistance to fixation. \\
N-tier architecture & Presentation, business, and data tiers have different responsibilities. \\
Client visibility and control & Scripts, markup, hidden fields, requests, cookies, and navigation are visible and mutable. \\
\bottomrule
\end{longtable}
\section{One-Sentence Takeaway}
\begin{CornellExamBox}{One-Sentence Takeaway}
\textbf{Web security requires the server to reestablish identity, authorization, state, canonical input, and output context on every request regardless of what the client interface appears to enforce.}
\end{CornellExamBox}
\end{document}
